\section[Задача ECDLP]{Задача ECDLP\ на эллиптических кривых}

Пусть $E(\mathbb{F}_p)$ --- эллиптическая кривая над простым конечным полем, а $\langle P\rangle=\{\mathcal{O},P,2P,\dots,(n-1)P\}$ --- конечная циклическая подгруппа порядка $n$. Если точка $Q$ принадлежит этой подгруппе, то существует единственное число $k \in \{0,1,\dots,n-1\}$, для которого выполняется равенство
\[
Q=kP.
\]
Задача дискретного логарифмирования на эллиптической кривой заключается в восстановлении числа $k$ по известным точкам $P$ и $Q$. Это число называется дискретным логарифмом точки $Q$ по основанию $P$ и обозначается $k=\log_PQ$.

Данная задача представляет интерес прежде всего потому, что она является естественной обратной задачей к скалярному умножению точки на число. Операция перехода от $k$ и $P$ к точке $Q=kP$ вычисляется эффективно (за $O(\log n)$) и лежит в основе большинства криптографических схем на эллиптических кривых. Обратный переход, то есть восстановление коэффициента $k$ по $P$ и $Q$, для корректно выбранных параметров считается вычислительно трудным. Именно на этой асимметрии основана практическая применимость эллиптических кривых в криптографии.

Смысл выбора базовой точки $P$ состоит в том, что после задания этой точки вся дальнейшая задача фактически рассматривается не на всей группе $E(\mathbb{F}_p)$, а внутри конкретной циклической подгруппы $\langle P\rangle$. Все её элементы имеют вид $kP$, где $k$ рассматривается по модулю $n$. Следовательно, открытые параметры (кривая, генератор $P$, точка $Q=kP$) определяют конечную циклическую группу известного порядка, а секретный параметр ($k$) интерпретируется как неизвестный коэффициент в представлении точки $Q$ через генератор $P$.

С вычислительной точки зрения задача дискретного логарифмирования на эллиптических кривых существенно отличается от аналогичной задачи в мультипликативных группах конечных полей. Для общего случая не известно алгоритмов субэкспоненциальной сложности, аналогичных методам индекс-калькуляса, успешно применяемым к обычной задаче дискретного логарифмирования в $\mathbb{F}_p^*$. Поэтому при корректном выборе параметров наиболее эффективными остаются общие алгоритмы вероятностного типа, трудоёмкость которых имеет порядок $O(\sqrt{n})$ групповых операций.